\chapter{Аналитическая часть}

\section{Синтетические данные}

Синтетические данные --– это сгенерированные компьютером данные,
обладающие теми же характеристиками и статистическими свойствами, что и
реальные данные.
Они применяются в том случае, если получение реальных данных затруднено, требует больших затрат, ставит под угрозу конфиденциальность пользователей или занимает слишком много времени.
Синтетические данные могут быть текстовыми или табличными, а
также представлять собой видео, изображения или звуки, как и реальные данные.~\cite{t}

Математически синтетические данные можно определить как выборку $X_{syn} = \{x_1, x_2, ..., x_n\}$ из некоторого распределения $P_{syn}$, которое аппроксимирует исходное распределение реальных данных $P_{real}$. Цель генерации синтетических данных заключается в минимизации расхождения между этими распределениями, которую описывают формулой~(\ref{eq:sint1}):

\begin{equation}\label{eq:sint1}
    \min_{P_{syn}} D(P_{real} || P_{syn}),
\end{equation}
где:
\begin{enumerate}[label*=---]
	\item  $D(P_{real} || P_{syn})$ --- мера расхождения между распределениями.
\end{enumerate}


Важным аспектом синтетических данных является их способность сохранять статистические свойства исходных данных.~\cite{sin2}

Это достигается через моделирование совместного распределения всех признаков по формуле~(\ref{eq:sint2}):

\begin{equation}\label{eq:sint2}
    P(X_1, X_2, ..., X_d) = \prod_{i=1}^{d} P(X_i | X_1, ..., X_{i-1}),
\end{equation}
где:
\begin{enumerate}[label*=---]
	\item $d$ --- размерность данных, где каждый условный множитель может быть смоделирован различными методами.
\end{enumerate}
Синтетические данные разделяют на три категории:
\begin{enumerate}[label*=---]
	\item Полностью синтетические данные, которые не имеют с реальными
	данными почти ничего общего и генерируются на основе некоторых признаков
	реальных данных. Они сохраняют статистические свойства реальных данных и
	фактически заменяют их.~\cite{t, sin2}
	\item Частично синтетические данные, которые заменяют только
	некоторые значения в реальных данных, если они являются конфиденциальными
	или содержат пропуски.~\cite{t, sin2}
	\item Гибридные синтетические данные, которые генерируются с
	использованием и реальных, и полностью синтетических данных: для реальных
	данных выбираются близкие им синтетические, и они объединяются, формируя
	гибридные данные.~\cite{t, sin2}
\end{enumerate}

% \subsubsection{Применение синтетических данных}

% Синтетические данные находят широкое применение в различных областях современной науки и технологий. В медицине они используются для обучения диагностических систем искусственного интеллекта без нарушения конфиденциальности пациентов. Финансовый сектор применяет синтетические данные для тестирования алгоритмов обнаружения мошенничества и оценки кредитных рисков.

%В области компьютерного зрения синтетические изображения позволяют создавать обширные обучающие выборки для нейронных сетей, особенно в условиях, когда сбор реальных данных затруднен или опасен. Автономные транспортные средства используют синтетические сценарии вождения для безопасного тестирования алгоритмов принятия решений.

%Математическая формализация качества синтетических данных может быть выражена через несколько критериев:

%1. \textbf{Точность (Fidelity):} мера того, насколько хорошо синтетические данные воспроизводят статистические свойства исходных данных:
%\begin{equation}
    %F = 1 - \frac{1}{m} \sum_{i=1}^{m} |S_{real}^{(i)} - S_{syn}^{(i)}|
%\end{equation}
%где $S_{real}^{(i)}$ и $S_{syn}^{(i)}$ --- $i$-я статистика реальных и синтетических данных соответственно.

%2. \textbf{Полезность (Utility):} способность синтетических данных обеспечивать сопоставимые результаты при решении задач машинного обучения:
%\begin{equation}
    %U = 1 - |Acc_{real} - Acc_{syn}|
%\end{equation}
%где $Acc_{real}$ и $Acc_{syn}$ --- точность модели, обученной на реальных и синтетических данных соответственно.

%3. \textbf{Приватность (Privacy):} мера защиты индивидуальной информации, часто выражаемая через расстояние до ближайшего реального примера:
%\begin{equation}
    %P = \min_{x \in X_{real}} \|x_{syn} - x\|_2
%\end{equation}

% \subsubsection{Проблемы и вызовы}

% Основные проблемы генерации синтетических данных включают в себя проблему режимного коллапса (mode collapse), когда генеративная модель создает данные только из ограниченного подмножества исходного распределения. Это особенно критично для высокоразмерных данных, где сложность распределения экспоненциально возрастает с увеличением размерности.

% Другой важной проблемой является балансирование между реалистичностью и приватностью. Чем более реалистичными являются синтетические данные, тем выше риск раскрытия конфиденциальной информации об исходных данных. Эта проблема математически формулируется как оптимизационная задача с ограничениями:

% \begin{equation}
%    \max_{P_{syn}} U(P_{syn}) \quad \text{subject to} \quad \mathcal{P}(P_{syn}) \geq \mathcal{P}_{min}
% \end{equation}

% где $U(P_{syn})$ --- функция полезности, а $\mathcal{P}(P_{syn})$ --- мера приватности с минимальным допустимым уровнем $\mathcal{P}_{min}$.

\subsubsection{Категориальные признаки}

Важную роль при генерации синтетических данных играют категориальные признаки, которые являются переменными, принимающими ограниченное число значений (например, пол, страна, тип товара).
Категориальные признаки позволяют выразить свойства и разбить данные на группы.~\cite{uch}

Математически категориальная переменная $X$ может быть представлена как случайная величина, принимающая значения из конечного множества ${C} = \{c_1, c_2, ..., c_k\}$ с вероятностями $P(X = c_i) = p_i$, где $\sum_{i=1}^{k} p_i = 1$.

Для многомерного случая с $m$ категориальными признаками совместное распределение описывается по формуле~(\ref{eq:11}):
\begin{equation}\label{eq:11}
    P(X_1 = c_{1,i_1}, X_2 = c_{2,i_2}, ..., X_m = c_{m,i_m}) = p_{i_1,i_2,...,i_m},
\end{equation}
где: 
\begin{enumerate}[label*=---]
	\item $p_{i_1,i_2,...,i_m}$ --- совместная вероятность соответствующей комбинации категорий.
\end{enumerate}
Категориальные признаки разделяют на три типа:
\begin{enumerate}[label*=---]
	\item Дискретные --- численные значения из конечного или счётного множества (например, количество комнат в квартире: $\{1, 2, 3, \dots\}$)
	\item Номинативные --- нет порядка (например, цвет: красный, синий, зелёный);
	\item Порядковые --- есть иерархия (например, уровень образования: школа, бакалавр, магистр);
	\item Бинарные --- два значения (например, да/нет).
\end{enumerate}

%От типа зависит выбор метода кодирования.
%Методы кодирования разбивают на:
%\begin{enumerate}[label*=---]
	%\item  Метка чисел (Label Encoding) --- каждой уникальной категории присваивается целое число от 0 до $N-1$, где $N$ - количество категорий. Математически это отображение $f: \mathcal{C} \rightarrow \{0, 1, ..., N-1\}$;
	%\item Среднее значение (Mean Encoding) --- замена категории средним значением целевой переменной для этой категории:
	%\begin{equation}
	    %\text{TE}(c_i) = \mathbb{E}[Y | X = c_i] = \frac{\sum_{j: x_j = c_i} y_j}{\sum_{j: x_j = c_i} 1}
	%\end{equation}
	%\begin{equation}
	    %\text{FE}(c_i) = \frac{\text{count}(X = c_i)}{n}
	%\end{equation}
	%где $n$ --- общее количество наблюдений;
	%\item Векторные представления (Embedding) --- преобразование категорий в плотные векторы фиксированной размерности $d$ через обучение: $\text{Emb}: \mathcal{C} \rightarrow \mathbb{R}^d$.
%\end{enumerate}

\subsubsection{Корреляция}

Корреляция в контексте генерации синтетических данных --- это статистическая взаимосвязь между признаками, которую необходимо сохранить, чтобы искусственные данные оставались реалистичными и пригодными для анализа или машинного обучения.
Если синтетические данные игнорируют связи, то неизбежны нереалистичные комбинации, снижение точности моделей, ошибки в анализах~\cite{uch}.

Корреляция между двумя числовыми переменными $X$ и $Y$ определяется как~(\ref{eq:1}):
\begin{equation}\label{eq:1}
    \rho_{X,Y} = \frac{\text{cov}(X,Y)}{\sigma_X \sigma_Y},
\end{equation}
где: 
\begin{enumerate}[label*=---]
	\item  $\sigma_X$, $\sigma_Y$ --- стандартные отклонения.
\end{enumerate}
%Проверить корреляцию можно несколькими способами:
%\begin{enumerate}[label*=---]
	%\item визуализация (корреляционные матрицы, диаграммы рассеяния);
	%\item статистические тесты (тест Спирмена);
	%\item метрики (коэффициент детерминации $R^2$, средняя абсолютная ошибка корреляций).
%\end{enumerate}

\subsubsection{Фробениусова норма}
Фробениусова норма --- это ключевая метрика в линейной алгебре и машинном обучении, используемая для оценки размера матриц и ошибок между ними.
В контексте генерации синтетических данных она измеряет, насколько синтетические данные отклоняются от реальных.

Для матрицы \( A \in \mathbb{R}^{m \times n} \) Фробениусова норма высчитывается по формуле~(\ref{eq:froben1}):

\begin{equation}\label{eq:froben1}
	\|A\|_F = \sqrt{\sum_{i=1}^{m}\sum_{j=1}^{n} |a_{ij}|^2},
\end{equation}
где:
\begin{enumerate}[label*=---]
	\item  \( a_{ij} \) --- элементы матрицы.
\end{enumerate}
\newpage
Для оценки качества синтетических данных, сгенерированных на основе исходных данных используется формула~(\ref{eq:froben2}):

\begin{equation}\label{eq:froben2}
	\varepsilon = \|P - \hat{P}\|_F = \sqrt{\sum_{i,j}(p_{ij} - \hat{p}_{ij})^2},
\end{equation}
где:
\begin{enumerate}[label*=---]
	\item \( P \) --- матрица распределения исходных данных;
	\item \( \hat{P} \) --- матрица синтетических данных.
	
\end{enumerate}


% \subsection{Оценка качества синтетических данных}

% Качество синтетических данных оценивается по нескольким критериям, каждый из которых имеет свои математические метрики.

% \subsubsection{Статистическое сходство}

% Основные статистические метрики включают:

%1. \textbf{Колмогоров-Смирнов тест} для сравнения одномерных распределений:
%\begin{equation}
    %D_{KS} = \sup_x |F_{real}(x) - F_{syn}(x)|
%\end{equation}
%где $F_{real}$ и $F_{syn}$ --- эмпирические функции распределения.

%2. \textbf{Wasserstein distance} (Earth Mover's Distance):
%\begin{equation}
    %W_p(\mu, \nu) = \left(\inf_{\gamma \in \Pi(\mu,\nu)} \int_{\mathcal{X} \times \mathcal{X}} d(x,y)^p d\gamma(x,y)\right)^{1/p}
%\end{equation}

%3. \textbf{Maximum Mean Discrepancy (MMD):}
%\begin{equation}
    %\text{MMD}^2(P, Q) = \|\mu_P - \mu_Q\|_{\mathcal{H}}^2
%\end{equation}
%где $\mu_P$, $\mu_Q$ --- средние в пространстве воспроизводящих ядер.

% \subsubsection{Машинное обучение метрики}

% Эффективность синтетических данных для задач машинного обучения оценивается через:

% 1. \textbf{Train on Synthetic, Test on Real (TSTR):}
%\begin{equation}
   % \text{TSTR} = \frac{\text{Acc}(\text{model trained on syn, tested on real})}{\text{Acc}(\text{model trained on real, tested on real})}
%\end{equation}

%2. \textbf{Train on Real, Test on Synthetic (TRTS):}
%\begin{equation}
    %\text{TRTS} = \frac{\text{Acc}(\text{model trained on real, tested on syn})}{\text{Acc}(\text{model trained on real, tested on real})}
%\end{equation}

\section{Методы генерации синтетических данных}

Создание синтетических данных включает в себя различные методы и подходы, которые позволяют генерировать данные, максимально приближенные к реальным по своим характеристикам и структуре~\cite{gan}.

%\subsubsection{Параметрические методы}

%Параметрические методы предполагают, что данные следуют известному распределению с набором параметров $\theta$. Процесс генерации заключается в оценке параметров $\hat{\theta}$ из исходных данных и последующей генерации новых выборок из распределения $p(x|\hat{\theta})$.

%Для нормального распределения:
%\begin{equation}
    %p(x|\mu,\sigma^2) = \frac{1}{\sqrt{2\pi\sigma^2}} \exp\left(-\frac{(x-\mu)^2}{2\sigma^2}\right)
%\end{equation}

%Оценки параметров методом максимального правдоподобия:
%\begin{align}
    %\hat{\mu} &= \frac{1}{n} \sum_{i=1}^{n} x_i \\
    %\hat{\sigma}^2 &= \frac{1}{n} \sum_{i=1}^{n} (x_i - \hat{\mu})^2
%\end{align}

%Для многомерного случая используется многомерное нормальное распределение:
%\begin{equation}
    %p(\mathbf{x}|{\mu},{\Sigma}) = \frac{1}{(2\pi)^{d/2}|l{\Sigma}|^{1/2}} \exp\left(-\frac{1}{2}(\mathbf{x}-{\mu})^T{\Sigma}^{-1}(\mathbf{x}-{\mu})\right)
%\end{equation}

%где $\mu$ --- вектор средних, ${\Sigma}$ --- ковариационная матрица .

%\subsubsection{Непараметрические методы}

%Непараметрические методы не делают предположений о форме распределения данных. Основные подходы включают:

%\begin{enumerate}[label*=---]

%\item Ядерная оценка плотности (KDE):
%\begin{equation}
    %\hat{f}(x) = \frac{1}{nh} \sum_{i=1}^{n} K\left(\frac{x-x_i}{h}\right)
%\end{equation}
%где $K$ --- ядерная функция, $h$ --- ширина окна.

%\item Гистограммный метод: разбиение пространства признаков на bins и моделирование дискретного распределения.

%\item Bootstrap-методы: повторная выборка с возвращением из исходных данных.
%\end{enumerate}
\subsection{GAN}

Генеративные состязательные нейросети --- это класс алгоритмов
искусственного интеллекта, используемых в машинном обучении без
вмешательства человека, реализованных системой двух нейронных сетей,
конкурирующих друг с другом в игровой среде с нулевой суммой.~\cite{gan}

Математическая основа GAN заключается в решении задачи минимакса. Пусть $p_{data}(x)$ --- распределение реальных данных, $p_z(z)$ --- распределение шума, $G(z;\theta_g)$ --- генератор с параметрами $\theta_g$, $D(x;\theta_d)$ --- дискриминатор с параметрами $\theta_d$.

Тогда целевую функцию GAN можно определить формулой~(\ref{eq:ganf1}):
\begin{equation}\label{eq:ganf1}
    \min_G \max_D V(D,G) = \mathbb{E}_{x \sim p_{data}(x)}[\ln D(x)] + \mathbb{E}_{z \sim p_z(z)}[\ln(1-D(G(z)))],
\end{equation}

%Процесс обучения включает альтернирующую оптимизацию:

%1. Обучение дискриминатора:
%\begin{equation}
    %\theta_d^{(t+1)} = \theta_d^{(t)} + \alpha \nabla_{\theta_d} \left[\frac{1}{m}\sum_{i=1}^{m} \log D(x^{(i)}) + %\frac{1}{m}\sum_{i=1}^{m} \log(1-D(G(z^{(i)})))\right]
%\end{equation}

%2. Обучение генератора:
%\begin{equation}
    %\theta_g^{(t+1)} = \theta_g^{(t)} - \alpha \nabla_{\theta_g} \left[\frac{1}{m}\sum_{i=1}^{m} \log(1-D(G(z^{(i)})))\right]
%\end{equation}

% где $\alpha$ --- скорость обучения, $m$ --- размер батча.


В системе GAN одна генеративная сеть $G$ генерирует
синтетические образцы, а другая сеть --- дискриминативная $D$, старается отличить правильные («оригинальные») образцы от неправильных (синтетических).~\cite{gan_r}

% Теоретический анализ GAN показывает, что при оптимальном дискриминаторе $D^*$ целевая функция минимизируется, когда $p_g = p_{data}$, где $p_g$ --- распределение, индуцированное генератором.

% Оптимальный дискриминатор:
%\begin{equation}
    %D^*(x) = \frac{p_{data}(x)}{p_{data}(x) + p_g(x)},
%\end{equation}

%При подстановке в целевую функцию получаем:
%\begin{equation}
    %C(G) = -\log(4) + 2 \cdot JSD(p_{data} || p_g),
%\end{equation}
%где $JSD$ --- дивергенция Дженсена-Шеннона:
%\begin{equation}
    %JSD(P||Q) = \frac{1}{2}KL(P||M) + \frac{1}{2}KL(Q||M),
%\end{equation}
%где $M = \frac{1}{2}(P+Q)$.
\newpage
Используя шум на входе, генеративная сеть $G$ создает новый образец, смешав несколько исходных образцов. Цель генератора можно описать выражением~(\ref{eq:gan1})
\begin{equation}\label{eq:gan1}
	\min _G \ln\{1-D[G(z)]\},
\end{equation}
где:
\begin{enumerate}[label*=---]
	\item $D$ --- функция дискриминатора;
	\item  $G$ --- функция генератора;
	\item  $z$ --- входные параметры.
\end{enumerate}


Дискриминативная сеть $D$ обучается различать подлинные и поддельные образцы, а результаты различения подаются на вход генеративной сети так, чтобы она смогла подобрать лучший набор параметров и дискриминативная сеть уже не смогла бы отличить подлинные образцы от поддельных.~\cite{gan_r}

При этом сеть $D$, анализируя образцы из оригинальных и синтетических данных, достигает некоторой точности различения.
Постепенно синтетические образцы на выходе генеративной сети становятся все более качественными. Цель дискриминатора $D$ можно описать формулой~(\ref{eq:gan2}):
\begin{equation}\label{eq:gan2}
	\max _D \{\ln D(x_t) + \ln (1-D[x_f])\},
\end{equation}
где:
\begin{enumerate}[label*=---]
	\item $x_t$ --- исходные данные; 
	\item $x_f$ --- сгенерированный образец.
\end{enumerate}



Таким образом, алгоритм представляет собой минимаксную состязание для двух
нейронных сетей, которое можно описать формулой~(\ref{eq:gan3}):
\begin{equation}\label{eq:gan3}
	\min _G \max _D \{\ln D(x_t) + \ln (1-D[G(z)])\},
\end{equation}

Для синтезации данных необходимо данные от генератора не только подавать на дискриминатор, но и складывать их в отдельную базу, особенно в фазе обучения, близкой к достижению минимаксного состояния равновесия.
В этом случае получится множество данных, очень близких к оригинальным, которые и являются искомыми синтетическими данными.~\cite{gan_r}

Структура сети GAN представлена на рисунке~\ref{img:gan1}
\includeimage
{gan1} % Имя файла без расширения (файл должен быть расположен в директории inc/img/)
{f} % Обтекание (с обтеканием)
{h!} % Положение рисунка (см. wrapfigure из пакета wrapfig)
{\textwidth} % Ширина рисунка
{Структура сети GAN} % Подпись рисунка

Структурная схема GAN, модифицированная для генерации синтетических данных, представлена на рисунке~\ref{img:gan2}.

\includeimage
{gan2} % Имя файла без расширения (файл должен быть расположен в директории inc/img/)
{f} % Обтекание (с обтеканием)
{h!} % Положение рисунка (см. wrapfigure из пакета wrapfig)
{\textwidth} % Ширина рисунка
{Структурная схема GAN, модифицированная для генерации синтетических данных} % Подпись рисунка

\clearpage
\subsection{cGAN}

Основное отличие cGAN (conditional GAN) от классической GAN заключается в том, что на дискриминатор и генератор в качестве дополнительного входного слоя подается дополнительный параметр $y$, который может являться условием при генерации.~\cite{cgan}

Математически cGAN можно представить формулой~(\ref{eq:cGAN1}):
\begin{equation}\label{eq:cGAN1}
    \min_G \max_D V(D,G) = \mathbb{E}_{x \sim p_{data}(x)}[\ln D(x|y)] + \mathbb{E}_{z \sim p_z(z)}[\ln(1-D(G(z|y)))],
\end{equation}
где:
\begin{enumerate}[label*=---]
	\item $y$ --- условная информация (метки классов, атрибуты).
\end{enumerate}
Генератор и дискриминатор принимают условную информацию различными способами:
\begin{enumerate}[label*=---]
	\item Конкатенация --- условие просто добавляется к входному вектору;
	\item Условная нормализация --- параметры нормализации зависят от условия;
	\item Attention механизмы --- условие влияет на веса внимания.
\end{enumerate}

На вход генератора подаются случайный шум и информация, которые объединяются в совместное входные параметры, в результате чего состязательная система обучения позволяет проявлять большую гибкость.
\newpage
Структурная схема cGAN представлена на рисунке~\ref{img:cgan}.
\includeimage
{cgan} % Имя файла без расширения (файл должен быть расположен в директории inc/img/)
{f} % Обтекание (с обтеканием)
{h!} % Положение рисунка (см. wrapfigure из пакета wrapfig)
{\textwidth} % Ширина рисунка
{Структурная схема cGAN } %

\subsection{ctGAN}

В традиционных GAN генератор повторяет распределение реальных данных.
В случае дисбаланса в категориальных колонках, если при обучении данные выбираются случайным образом, то строки не будут представлены в достаточной степени и, следовательно, генератор может быть обучен неправильно.
Для решения этих проблем была предложена архитектура cGAN, адаптированная для генерации равномерно распределенных многомодальных табличных данных, названная ctGAN (conditional tabular GAN).~\cite{ctgan}

%Ключевые инновации ctGAN:

%\begin{enumerate}[label*=---]

%\item Mode-specific normalization --- для каждого непрерывного столбца используется смесь гауссианов:
%\begin{equation}
	%p(x) = \sum_{i=1}^{K} \pi_i \mathcal{N}(x; \mu_i, \sigma_i^2)
%\end{equation}

%где параметры $\{\pi_i, \mu_i, \sigma_i\}$ оцениваются методом EM-алгоритма.

%\item Conditional sampling --- во время обучения происходит условная выборка на основе дискретных столбцов для %обеспечения равномерного представления всех категорий.

	%\item Training-by-sampling --- вероятность выбора строки для обучения:
%\begin{equation}
	%p(\text{row } i) = \frac{1}{|\text{rows with same discrete values as row } i|}
%\end{equation}
%\end{enumerate}


Задача генератора ctGAN состоит в том, чтобы все категории из дискретных атрибутов были равномерно отобраны в процессе обучения, а во время тестирования было восстановлено реальное распределение данных.~\cite{kim2024ctgan}
%В основу идей CTGAN положено следующее соображение: пусть $k*$ --- значение из $і*-го$ дискретного столбца $D_{i*}$, которому должны соответствовать сгенерированные выборки $r$.
%Тогда генератор можно интерпретировать как условное распределение строк с учетом конкретного значения в данном столбце.

\newpage
Алгоритм ctGAN включает следующие этапы:
\begin{enumerate}
	\item Предобработка данных: анализ каждого столбца для определения типа (непрерывный/дискретный);

	\item Условная выборка: выбор дискретного столбца $col_i$ с вероятностью пропорциональной логарифму числа уникальных значений, выбор конкретного значения $val_j$ из столбца равномерно, выборка строк, содержащих выбранное значение;

	\item Генерация условного вектора: создание маски $m$ для указания условного столбца, формирование вектора условий $c$.
\end{enumerate}
Математически процесс генерации можно записать по формуле~(\ref{eq:ctGAN1}):
\begin{equation}\label{eq:ctGAN1}
    {x}_{syn} = G({z}, {c}, {m}),
\end{equation}
где:
\begin{enumerate}[label*=---]
	\item ${z} \sim \mathcal{N}(0, {I})$ --- случайный шум;
	\item $c$ --- условный вектор;
	\item $m$ --- маска.
\end{enumerate}


\subsection{VAE}

Вариационные автокодировщики (VAE) представляют собой генеративные модели, основанные на принципах вариационного вывода и теории информации. VAE состоит из двух основных компонентов: кодировщика (encoder) и декодировщика (decoder).~\cite{vae}

Математическая формулировка VAE основана на максимизации нижней границы логарифма правдоподобия данных и представляется выражением~(\ref{eq:vae1}):

\begin{equation}\label{eq:vae1}
    \ln p({x}) \geqslant \mathbb{E}_{q_\phi({z}|{x})}[\ln p_\theta({x}|{z})] - KL(q_\phi({z}|{x}) || p({z})),
\end{equation}
где:
\begin{enumerate}[label*=---]
	\item $q_\phi({z}|{x})$ --- приближенное апостериорное распределение (кодировщик);
	\item  $p_\theta({x}|{z})$ --- декодер;
	\item  $p({z})$ --- априорное распределение (обычно стандартное нормальное);
	\item  $KL$ --- дивергенция Кульбака-Лейблера.
\end{enumerate}
%Кодировщик параметризуется как:
%\begin{align}
    %\mu_\phi({x}) &= f_\mu({x}), \\
    %\sigma_\phi({x}) &= f_\sigma({x}),\\
   % q_\phi({z}|{x}) &= \mathcal{N}({z}; \mu_\phi({x}), \sigma_\phi^2({x})),
%\end{align}

%Трюк репараметризации позволяет сделать процесс дифференцируемым:
%\begin{equation}
    %{z} = \mu_\phi({x}) + \sigma_\phi({x}) \odot {\varepsilon}
%\end{equation}
%где ${\varepsilon} \sim \mathcal{N}(0, {I})$.

% \subsection{Нормализующие потоки (Normalizing Flows)}

% Нормализующие потоки представляют собой последовательность обратимых преобразований, которые преобразуют простое распределение (например, стандартное нормальное) в сложное целевое распределение.

% Пусть $f_1, f_2, ..., f_K$ --- последовательность обратимых функций. Тогда:
%\begin{equation}
    %\mathbf{z}_K = f_K \circ f_{K-1} \circ ... \circ f_1(\mathbf{z}_0)
%\end{equation}

%Плотность преобразованного распределения вычисляется по формуле замены переменных:
%\begin{equation}
    %\log p(\mathbf{z}_K) = \log p(\mathbf{z}_0) - \sum_{k=1}^{K} \log \left| \det \frac{\partial f_k}{\partial \mathbf{z}_{k-1}} \right|
%\end{equation}

%Популярные типы преобразований:
%1. \textbf{Planar flows:} $f(\mathbf{z}) = \mathbf{z} + \mathbf{u} \tanh(\mathbf{w}^T \mathbf{z} + b)$
%2. \textbf{Radial flows:} $f(\mathbf{z}) = \mathbf{z} + \beta h(\alpha, r)(\mathbf{z} - \mathbf{z}_0)$
%3. \textbf{Coupling layers:} разделение вектора на две части с преобразованием только одной части


\subsection{Байесовские сети}

Статистические и вероятностные методы представляют собой фундаментальный подход к генерации синтетических данных, основанный на математическом моделировании распределений реальных данных.~\cite{uch}

Основные принципы статистических методов:
\begin{enumerate}[label*=---=]
	\item оценка параметров распределения --- по исходным данным оцениваются параметры теоретического распределения;
	\item валидация модели --- проверка соответствия выбранной модели данным;
	\item генерация новых выборок --- использование оцененных параметров для создания синтетических данных.
\end{enumerate}


Для многомерных данных часто используется предположение о независимости признаков:
\begin{equation}
	p(\mathbf{x}) = \prod_{i=1}^{d} p(x_i),
\end{equation}


Байесовские сети --- это тип вероятностной графической модели,
которая представляет отношения между переменными с использованием
ориентированного ациклического графа.~\cite{uch}


В Байесовских сетях:
\begin{enumerate}
	\item узлы соответствуют случайным переменным;
	\item дуги отражают причинно-следственные связи;
	\item каждому узлу сопоставлено условное распределение вероятностей.
\end{enumerate}

\newpage
Математическая основа Байесовских сетей базируется на разложении совместного распределения по формуле~(\ref{eq:baes1}):
\begin{equation}\label{eq:baes1}
    P(X_1, X_2, ..., X_n) = \prod_{i=1}^{n} P(X_i | Pa(X_i)),
\end{equation}
где:
\begin{enumerate}[label*=---]
	\item $Pa(X_i)$ --- родительские узлы для узла $X_i$.
\end{enumerate}


Процесс обучения байесовской сети включает структурное обучение: определение графа зависимостей, поиск на основе оценок, тестирование условной независимости, жадные алгоритмы и эволюционные подходы; параметрическое обучение: оценка условных распределений, максимальное правдоподобие, байесовский подход с априорными распределениями~\cite{bs}.

Для генерации синтетических данных используется прямая выборка~(\ref{eq:baes2}):
\begin{equation}\label{eq:baes2}
    x_i \sim P(X_i | pa(x_i)),
\end{equation}
где:
\begin{enumerate}[label*=---]
	\item $pa(x_i)$ --- значения родительских узлов.
\end{enumerate}


% \subsection{Методы на основе деревьев решений}

% Деревья решений и их ансамбли также могут использоваться для генерации синтетических данных. Основная идея заключается в построении модели, которая описывает структуру данных через иерархическую систему правил.

% \subsubsection{CART для генерации данных}

% Алгоритм CART (Classification and Regression Trees) может быть адаптирован для генерации:

% 1. \textbf{Построение дерева:} на основе исходных данных строится дерево решений
% 2. \textbf{Анализ листьев:} в каждом листе анализируется распределение целевых значений
% 3. \textbf{Генерация:} новые данные создаются путем прохода по дереву и выборки из листовых распределений

% Критерий разбиения для непрерывных переменных:
% \begin{equation}
    %\text{Gini}(S) = 1 - \sum_{i=1}^{c} p_i^2
%\end{equation}

%где $p_i$ --- доля класса $i$ в множестве $S$.

% \subsubsection{Random Forest для синтеза}

% Случайный лес может использоваться для более робустной генерации:
% \begin{equation}
    %\hat{y} = \frac{1}{B} \sum_{b=1}^{B} T_b(\mathbf{x})
%\end{equation}

% где $T_b$ --- $b$-е дерево, $B$ --- количество деревьев.

\subsection{Сравнение методов}

Cравнение методов генерации синтетических данных представлено в таблице~\ref{tab:my_table_label}:

\begin{itemize}[label*=---]
	\item 1 --- тип данных. Возможность метода те или иные типы данных;
	\item 2 --- сложность обучения, где $E$ --- эпохи; $N$ --- слои; $B$ --- батч; $Z$ --- число узлов; $D$ --- количество зависимостей;
	% \item 3 --- уровень дифференциальной приватности. Наличие у метода функционала дифференциальной приватности;
	\item 3 --- качество сохранения корреляций между признаками по Фробениусовой норме;
	% \item 4 --- способность работать с пропущенными значениями;
	\item 4 --- модификация для ДП. Возможности модификации алгоритма для обеспечения ДП.
\end{itemize}

\newpage
\begin{table}[ht]
	\centering
	\captionsetup{justification=raggedleft}
	\caption{Cравнение методов генерации синтетических данных}
	\label{tab:my_table_label}
\begin{tabularx}{\textwidth}{|X|X|X|X|X|}
	\hline
	Метод & 1 & 2 & 3 & 4\\
	\hline
	GAN & Изображение & $O(E \cdot N \cdot B)$  & 0.07 & DP-SGD при обучении дискриминатора \\
	\hline
	cGAN & Изображение & $O(E \cdot N \cdot B)$ & 0.07  &  DP-SGD для защиты приватности условных меток \\
	\hline
	ctGAN & Табличные данные & $O(E \cdot N \cdot B)$  & 0.06 & DP-SGD для добавления шума к трансформированным данным \\
	\hline
	VAE & Дискретные данные & $O(E \cdot N \cdot B)$ & 0.11 & DP-SGD при обучении энкодера и декодера \\
	\hline
	БС & Табличные данные & $O(Z \cdot D )$ & 0.07 & Добавление шума к частотным таблицам\\
	\hline
\end{tabularx}
\end{table}
\section{Дифференциальная приватность}

Дифференциальная приватность --- это математическая концепция, которая обеспечивает информационно-теоретическую гарантию конфиденциальности.
Стандартным способом реализации дифференциальной конфиденциальности является введение достаточного количества шума в данные.

Однако в дополнение к другим ограничениям дифференциальной конфиденциальности этот процесс добавления шума влияет на точность и полезность данных, но минимизирует влияние индивидуальных записей на результаты статистических запросов.~\cite{dp}

\subsection{Теоретические основы дифференциальной приватности}

Дифференциальная приватность базируется на принципе неразличимости: результаты анализа не должны позволять определить, участвовал ли конкретный индивид в исследовании или нет. Этот принцип формализуется через понятие соседних наборов данных.

Два набора данных $D$ и $D'$ называются соседними, если они отличаются записью ровно одного индивида $|D \triangle D'| = 1$, где $\triangle$ обозначает симметрическую разность множеств.

Дифференциальная приватность обеспечивает защиту личной информации, гарантируя, что результаты анализа не позволят точно определить информацию о конкретных индивидах. 
В рамках ДП данные обрабатываются с добавлением шума, что минимизирует влияние индивидуальных записей на результаты статистических запросов.

Таким образом, как показано на рисунке~\ref{img:diff}, даже если в набор данных добавляется или исключается информация о конкретном индивиде, это не влияет на общие выводы о группе.~\cite{Wood2018DifferentialPA}

%Представим БД, содержащую медицинские записи 100 человек, из которых 20 страдают диабетом.
%Злоумышленник хочет узнать, есть ли у человека $X$ диабет, и уже узнал, что 19 из остальных 99 человек в выборке являются диабетиками.
%Запрашивая базу данных и получая информацию о том, что в выборке 20 диабетиков, злоумышленник делает вывод, что человек
%$X$ также страдает диабетом.
%Этот процесс извлечения дедуктивной информации называется дифференцированной атакой.

%Теперь представим, что доступ к данным, хранящихся в БД, предоставляется через систему запросов с дифференциальной приватностью.
%Когда злоумышленник делает запрос, система возвращает истинное значение с некоторым случайным шумом.
%Даже если злоумышленник знает все о других 99 людях, то он не может быть уверен, есть ли у человека $X$ диабет.
%Конфиденциальность человека $X$ защищена, поскольку его риск нарушения конфиденциальности остается примерно одинаковым независимо от того, входит ли он в набор данных или нет.
%Это справедливо для любого отдельного человека и любого набора данных, независимо от того, насколько необычными или отличительными могут быть данные конкретного индивидуума.

\includeimage{diff}{f}{h}{1\textwidth}{Принцип дифференциальной приватности}

% \subsection{Глобальная и локальная дифференциальная приватность}

% \subsubsection{Глобальная дифференциальная приватность}

%В модели глобальной дифференциальной приватности существует доверенный куратор данных, который собирает все данные, применяет приватные механизмы и публикует результаты. Эта модель предполагает, что куратор является честным, но любопытным.

%Преимущества глобальной модели:
%- Высокая точность результатов
%- Эффективное использование приватного бюджета
%- Возможность сложных аналитических запросов

%Недостатки:
%- Требует доверия к куратору
%- Единая точка отказа
%- Сложности с распределенными данными

%\subsubsection{Локальная дифференциальная приватность}

% модели локальной дифференциальной приватности каждый индивид самостоятельно добавляет шум к своим данным перед их передачей.

%Алгоритм $A$ обеспечивает $\varepsilon$-локальную дифференциальную приватность, если для любых двух входных значений $x$ и $x'$ и любого множества выходов $S$:
%begin{equation}
   % \frac{Pr[A(x) \in S]}{Pr[A(x') \in S]} \leq e^\varepsilon
%\end{equation}

%Основные механизмы локальной дифференциальной приватности:

%1. \textbf{Randomized Response:} для бинарных данных
%\begin{equation}
    %Pr[A(x) = x] = \frac{e^\varepsilon}{1 + e^\varepsilon}
%\end{equation}
%\begin{equation}
    %Pr[A(x) = 1-x] = \frac{1}{1 + e^\varepsilon}
%\end{equation}

%2. \textbf{RAPPOR:} для категориальных данных
%3. \textbf{Duchi et al.:} для числовых данных в многомерном пространстве

% \subsection{Приватность при обучении машинного обучения}

% \subsubsection{Дифференциально-приватное обучение}

%Применение дифференциальной приватности к обучению моделей машинного обучения представляет особый интерес. Основные подходы:
%\begin{enumerate}[label*=---]
%\item Добавление шума к градиентам:
%\begin{equation}
	%\mathbf{g}_t = \frac{1}{|\mathcal{B}|} \sum_{i \in \mathcal{B}} \text{clip}(\nabla_\theta \ell(\theta, \mathbf{x}_i), C) + \mathcal{N}(0, \sigma^2 C^2 \mathbf{I})
%\end{equation}

%где $\text{clip}(\mathbf{g}, C) = \mathbf{g} \cdot \min(1, C/\|\mathbf{g}\|_2)$ --- функция обрезки градиента.

%\item Добавление шума к параметрам модели:
%\begin{equation}     
	%\tilde{\theta} = \theta + \mathcal{N}(0, \sigma^2 \mathbf{I})
%\end{equation}

%\item Приватная агрегация:
%\begin{equation}
	%\hat{f}(x) = \arg\max_y \left( \sum_{i=1}^{n} \mathbf{1}[f_i(x) = y] + \text{Lap}\left(\frac{1}{\varepsilon}\right) \right)
%\end{equation}
%\end{enumerate}


%\subsubsection{Анализ приватности для итеративных алгоритмов}

%Для анализа накопления приватного бюджета в итеративных алгоритмах используются продвинутые техники:

%1. \textbf{Moments Accountant:} отслеживание моментов приватной функции потерь
%2. \textbf{Rényi Differential Privacy:} обобщение классической дифференциальной приватности:
%\begin{equation}
    %D_\alpha(P||Q) = \frac{1}{\alpha-1} \log\left(\sum_x P(x)^\alpha Q(x)^{1-\alpha}\right)
%\end{equation}

%3. \textbf{Zero-Concentrated Differential Privacy:} использование концентрационных неравенств

\newpage
\section{Типы приватности}

Параметр приватности $\varepsilon$ --- это неотрицательное значение, определяющее уровень приватности, который влияет на количество добавляемого шума.~\cite{hu2015differential}

Интуитивная интерпретация параметра $\varepsilon$:
\begin{enumerate}[label*=---]
	\item $\varepsilon = 0$ --- идеальная приватность (бесконечный шум)
	\item $\varepsilon \to \infty$ --- отсутствие приватности
	\item $\varepsilon \in (0, 1]$ --- строгая приватность
	\item $\varepsilon \in (1, 10]$ --- умеренная приватность
	\item $\varepsilon > 10$ --- слабая приватность
\end{enumerate}


\subsection{$\varepsilon$ -- дифференцированная приватность}

Модель $\varepsilon$ -- дифференцированной приватности гарантирует, что результаты анализа данных не изменятся значительно, если одна запись будет добавлена или удалена.~\cite{hu2015differential}

Механизм $M$ обеспечивает $\varepsilon$ -- дифференциальную приватность, если для любых двух соседних наборов данных $D$ и $D'$ из множества $S$ выполняется условие по формуле~(\ref{equ:eps-diff}):

\begin{equation}\label{equ:eps-diff}
	Pr[M(D) \in S] \leqslant \exp(\varepsilon) \times Pr[M(D') \in S],
\end{equation}
где:
\begin{enumerate}[label*=---]
	\item $\varepsilon$ --- параметр приватности.
\end{enumerate}


%Эквивалентно, это можно записать через отношение плотностей:
%\begin{equation}
    %\left|\log\frac{Pr[M(D)]}{Pr[M(D')]}\right| \leq \varepsilon
%\end{equation}

\subsection{$(\varepsilon, \delta)$ -- дифференцированная приватность}

Модель $(\varepsilon, \delta)$ -- дифференцированной приватности вводит дополнительный параметр $\delta$, который ослабляет строгие гарантии приватности.
Параметр $\delta$ определяет вероятность того, что гарантия приватности будет нарушена в некоторых случаях.
Это дает гибкость в балансе между приватностью и точностью результатов.~\cite{hu2015differential}

\newpage
Механизм $M$ обеспечивает ($\varepsilon$, $\delta$) -- дифференциальную приватность, если для любых двух соседних наборов данных $D$ и $D'$ (различающихся на одну запись) выполняется условие по формуле~(\ref{equ:eps-delta-diff}):

\begin{equation}\label{equ:eps-delta-diff}
	Pr[M(D) \in S] \leqslant \exp(\varepsilon) \times Pr[M(D') \in S] + \delta,
\end{equation}
где:
\begin{enumerate}[label*=---]
	\item $\varepsilon$ --- параметр приватности;
	\item $\delta$ --- параметр вероятности нарушения приватности.
\end{enumerate}


%Параметр $\delta$ позволяет:
%- Использовать более эффективные механизмы
%- Работать с алгоритмами, требующими точных вычислений
%- Применять дифференциальную приватность к более широкому классу задач

%Типичные значения: $\delta \leq \frac{1}{n^2}$ или $\delta \leq \frac{1}{n\log n}$, где $n$ --- размер выборки.

%\subsection{Концентрированная дифференциальная приватность}

%Концентрированная дифференциальная приватность (CDP) представляет собой усиление стандартного определения, которое лучше подходит для анализа композиции механизмов.

%$\textbf{Определение (Zero-CDP):} Механизм $M$ удовлетворяет $\rho$-zero-concentrated differential privacy, если для всех соседних $D, D'$ и всех $\alpha \in (1, \infty)$:
%\begin{equation}
    %D_\alpha(M(D) || M(D')) \leq \rho \alpha
%\end{equation}

%где $D_\alpha$ --- $\alpha$-дивергенция Реньи.

%\subsection{Rényi дифференциальная приватность}

%RDP обобщает классическую дифференциальную приватность, используя дивергенцию Реньи:

%\textbf{Определение (RDP):} Механизм $M$ удовлетворяет $(\alpha, \varepsilon)$-RDP, если для всех соседних $D, D'$:
%\begin{equation}
    %D_\alpha(M(D) || M(D')) \leq \varepsilon
%\end{equation}

%где $D_\alpha(P || Q) = \frac{1}{\alpha-1} \log \mathbb{E}_{X \sim Q}\left[\left(\frac{P(X)}{Q(X)}\right)^\alpha\right]$ для $\alpha > 1$.

%Преимущества RDP:
%- Более точная оценка композиции
%- Лучшие границы для итеративных алгоритмов
%- Упрощенный анализ для гауссовского механизма

%\subsubsection{Практические соображения выбора параметров}

%Выбор параметров $\varepsilon$ и $\delta$ зависит от нескольких факторов:
%\begin{enumerate}[label*=---]

%\item размер выборки --- большие выборки позволяют использовать меньшие значения $\varepsilon$;
%\item чувствительность данных --- медицинские данные требуют строже приватности;
%\item требуемая точность --- высокая точность требует больших значений $\varepsilon$;
%\item количество запросов --- композиция требует распределения приватного бюджета.
%\end{enumerate}


%Рекомендуемые диапазоны:
%- Высокая приватность: $\varepsilon \in [0.1, 1.0]$, $\delta \leq 10^{-5}$
%- Умеренная приватность: $\varepsilon \in [1.0, 3.0]$, $\delta \leq 10^{-4}$
%- Низкая приватность: $\varepsilon \in [3.0, 10.0]$, $\delta \leq 10^{-3}$
\section{Механизмы приватности}
\subsection{Механизм Лапласа}

Лапласов механизм добавляет к результату запроса шум, сгенерированный из распределения Лапласа~\cite{db}.

%Распределение Лапласа с параметрами местоположения $\mu$ и масштаба $b$ имеет плотность, которая вычисляется по формуле~\ref{eq:lapl}:
%\begin{equation}\label{eq:lapl}
    %f(x|\mu,b) = \frac{1}{2b}\exp\left(-\frac{|x-\mu|}{b}\right),
%\end{equation}

%Для обеспечения $\varepsilon$--дифференциальной приватности параметр масштаба выбирается как $b = \frac{\Delta f}{\varepsilon}$.

Результат Лапласова механизма вычисляется по формуле~(\ref{eq:laplas}):
\begin{equation}\label{eq:laplas}
	M'(D) = f(D) + \text{Lap}\left(0, \frac{\Delta f}{\varepsilon}\right),
\end{equation}
где:
\begin{enumerate}[label*=---]
	\item $\Delta f$ --- глобальная чувствительность функции $f$;
	\item  $\varepsilon$ --- параметр приватности.
\end{enumerate}


%Свойства Лапласова механизма:
%- Несмещенность: $\mathbb{E}[M'(D)] = f(D)$
%- Дисперсия: $\text{Var}[M'(D)] = \frac{2(\Delta f)^2}{\varepsilon^2}$
%- Оптимальность для $L_1$-чувствительности

\subsection{Механизм Гаусса}

Гауссов механизм похож на Лапласов, но шум добавляется на основе нормального (Гауссова) распределения.~\cite{db}
К результатам запроса добавляется шум по формуле~(\ref{eq:gauss-norm}):

\begin{equation}\label{eq:gauss-norm}
	M'(D) = f(D) + \mathcal{N}(0, \sigma^2 ),
\end{equation}
где:
\begin{enumerate}[label*=---]
	\item $\mathcal{N}(0, \sigma^2)$ --- многомерное нормальное распределение.
\end{enumerate}


Для обеспечения $(\varepsilon, \delta)$-дифференциальной приватности параметр $\sigma$ вычисляется по формуле~(\ref{eq:gauss-sigma}):
\begin{equation}\label{eq:gauss-sigma}
	\sigma = \frac{\Delta_2 f \sqrt{2 \ln(1.25/\delta)}}{\varepsilon},
\end{equation}
где:
\begin{enumerate}[label*=---]
	\item $\Delta_2 f$ --- $L_2$-чувствительность функции.
\end{enumerate}

%\begin{equation}
    %\Delta_2 f = \max_{D \sim D'} \|f(D) - f(D')\|_2,
%\end{equation}

%Преимущества Гауссова механизма:
%- Лучше подходит для многомерных выходов
%- Более эффективен при малых значениях $\delta$
%- Естественная композиция через сложение дисперсий

\subsection{Экспоненциальный механизм}

Экспоненциальный механизм работает путем добавления случайности в
процесс выбора элементов, где вероятность выбора каждого элемента зависит от его качества (оценки полезности) и параметра приватности.~\cite{db}

Пусть $\mathcal{R}$ --- множество возможных выходов, $u: \mathcal{D} \times \mathcal{R} \to \mathbb{R}$ --- функция полезности. Экспоненциальный механизм выбирает выход $r \in \mathcal{R}$ с вероятностью пропорциональной:
\begin{equation}\label{eq:exp-mech}
	P[M(D)] \propto \exp\left(\frac{\varepsilon u(D, r)}{2\Delta u}\right),
\end{equation}
где:
\begin{enumerate}[label*=---]
	\item $\Delta u$ --- чувствительность функции полезности, вычислемая по формуле~(\ref{eq:exp}):
\end{enumerate}

\begin{equation}\label{eq:exp}
    \Delta u = \max_{r \in \mathcal{R}} \max_{D \sim D'} |u(D, r) - u(D', r)|,
\end{equation}
Нормализованная вероятность:
\begin{equation}
    P[M(D)] = \frac{\exp\left(\frac{\varepsilon u(D, r)}{2\Delta u}\right)}{\sum_{r' \in \mathcal{R}} \exp\left(\frac{\varepsilon u(D, r')}{2\Delta u}\right)},
\end{equation}

%Применения экспоненциального механизма:
%- Приватный выбор наилучшего параметра
%- Выпуск топ-k элементов
%- Аукционы с дифференциальной приватностью


\newpage
\subsection{Сравнение механизмов}

Сравнение механизмов дифференциальной приватности представлено в таблице~\ref{tab:dp}:

\begin{itemize}[label*=---]
	\item 1 --- тип приватности. Тип приватности, который использует данный метод;
	\item 2 --- тип данных. Типы данных, для которых подходит механизм;
	\item 3 --- вычислительная сложность. Сложность применения механизма.
\end{itemize}

\begin{table}[ht]
	\centering
	\captionsetup{justification=raggedleft}
	\caption{Сравнение механизмов дифференциальной приватности}
	\label{tab:dp}
\begin{tabularx}{\textwidth}{|X|X|X|X|}
	\hline
	Механизм & 1 & 2 & 3 \\
	\hline
	Лапласовский &  $(\varepsilon, 0)$ & Числовые & $O(1)$ \\
	\hline
	Гауссовский & $(\varepsilon, \delta)$ & Числовые & $O(1)$ \\
	\hline
	Экспоненциальнй& $\varepsilon$ & Категориальные & $O(m)$ \\
	\hline
\end{tabularx}
\end{table}
\section{Формализация задачи}

В ходе выполнения выпускной квалификационной работы необходимо разработать и реализовать метод, предоставляющий возможность генерировать синтетические данные с гарантиями дифференциальной приватности.
Концепция метода заключается в совместном использовании метода генерации синтетических данных и метода дифференциальной приватности путем добавления шума на этапе генерации.

\subsection{Идея комбинированного метода}

Пусть $D = \{x_1, x_2, ..., x_n\}$ --- исходная база данных, где каждая запись $x_i \in \mathbb{R}^d$ представляет собой $d$-мерный вектор признаков. Задача состоит в разработке метода $\mathcal{M}$, который:
\begin{enumerate}[label*=---]
\item генерирует синтетические данные: $D_{syn} = \mathcal{M}(D) = \{y_1, y_2, ..., y_m\}$, где $m$ --- заданное количество синтетических записей;
\item обеспечивает дифференциальную приватность: $\mathcal{M}$ удовлетворяет $(\varepsilon, \delta)$-дифференциальной приватности;
\item сохраняет полезность: синтетические данные сохраняют статистические свойства исходных данных.
\end{enumerate}


Также метод должен быть реализован как приложение с пользовательским интерфейсом.
На рисунке~\ref{img:idef0} представлена функциональная модель задачи в виде
диаграммы IDEF0 нулевого уровня.

\includeimage
{idef0} % Имя файла без расширения (файл должен быть расположен в директории inc/img/)
{f} % Обтекание (с обтеканием)
{h!} % Положение рисунка (см. wrapfigure из пакета wrapfig)
{\textwidth} % Ширина рисунка
{Функциональная модель задачи} %

%\section{Формат табличных данных}

%В рамках разрабатываемой системы поддерживаются следующие форматы табличных данных, каждый со своими особенностями и областями применения.

%\subsection{CSV}

%CSV (Comma-Separated Values) --- текстовый формат, предназначенный для представления табличных данных. 
%Строка таблицы соответствует строке текста, которая содержит одно или несколько полей, разделенных запятыми.
%В первой строке допускается указать наименования столбцов данных, но это не является обязательным требованием.
%Данный формат не предполагает типизацию данных, все значения данных считаются строками.
%CSV является наиболее часто используемым форматом для хранения структурированных данных.

%Математическое представление CSV-файла:
%%CSV = \{H, R\}
%\end{equation}
%где $H = \{h_1, h_2, ..., h_d\}$ --- заголовки столбцов, $R = \{r_1, r_2, ..., r_n\}$ --- строки данных, каждая $r_i = \{r_{i,1}, r_{i,2}, ..., r_{i,d}\}$.

%Преимущества CSV:
%- Простота формата
%- Универсальная поддержка
%- Человекочитаемость
%- Малый размер файлов

%Недостатки:
%- Отсутствие типизации
%- Проблемы с кодировкой
%- Ограниченная поддержка сложных структур

%\subsection{Parquet}

%Parquet --- это открытый, колонно-ориентированный формат хранения данных, который используется для эффективной обработки больших объёмов данных в распределённых системах.
%Он был создан как part экосистемы Hadoop и поддерживает аналитику и обработку данных в таких распределённых системах как Apache Hadoop,
%Apache Spark, Apache Drill и других.

%Структура Parquet файла:
%\begin{equation}
    %\text{Parquet} = \{\text{Header}, \text{Row Groups}, \text{Footer}\}
%\end{equation}

%где каждая Row Group содержит Column Chunks:
%\begin{equation}
    %\text{Row Group} = \{\text{Column Chunk}_1, ..., \text{Column Chunk}_d\}
%\end{equation}

%Каждый Column Chunk включает:
%- Data Pages: сжатые данные столбца
%- Dictionary Pages: словарь для кодирования
%- Statistics: минимум, максимум, количество null

%Преимущества Parquet:
%- Эффективное сжатие (до 75% экономии места)
%- Быстрое чтение столбцов
%- Поддержка сложных типов данных
%- Встроенная статистика

%Недостатки:
%- Сложность формата
%- Неэффективность для построчного чтения
%- Требует специализированных библиотек

%\subsection{Avro}

%Формат Avro - это компактный, быстрый и сериализуемый формат данных, разработанный
%Apache Software Foundation. Он обеспечивает схему данных в JSON формате, что делает его самоописывающим. Avro поддерживает различные типы данных, сжатие данных и является эффективным для передачи данных между различными системами. Он широко используется в Big Data и распределенных системах для обмена данными и хранения данных.

%Схема Avro описывается в JSON:
%\begin{equation}
 %   \text{Schema} = \{\text{type}, \text{name}, \text{fields}\}
%\end{equation}

%где fields определяет структуру записи:
%\begin{equation}
%    \text{fields} = \{[\text{name}_1, \text{type}_1], ..., [\text{name}_d, \text{type}_d]\}
% \end{equation}

%Поддерживаемые типы данных:
%- Примитивные: null, boolean, int, long, float, double, bytes, string
%- Сложные: record, enum, array, map, union, fixed

%Преимущества Avro:
%- Компактная сериализация
%- Эволюция схемы
%- Богатые типы данных
%- Кросс-языковая поддержка

%Недостатки:
%- Сложность схемы
%- Overhead для малых файлов
%- Требует понимания схемы для чтения


% \end{tabularx}

%\section{Заключение аналитической части}

%В данной аналитической части была проведена комплексная оценка современного состояния методов генерации синтетических данных и техник обеспечения дифференциальной приватности.

%Основные результаты анализа:

%1. \textbf{Синтетические данные} представляют собой перспективное направление для решения проблем недостатка данных и обеспечения приватности. Рассмотренные методы показывают различную эффективность в зависимости от типа данных и требований к качеству.

%2. \textbf{Генеративно-состязательные сети} демонстрируют лучшие результаты для сложных распределений, но требуют значительных вычислительных ресурсов и тщательной настройки.

%3. \textbf{Дифференциальная приватность} обеспечивает математически строгие гарантии конфиденциальности, но требует баланса между приватностью и полезностью данных.

%4. \textbf{Комбинирование методов} генерации синтетических данных с техниками дифференциальной приватности представляет собой актуальную исследовательскую задачу с практической значимостью.

%Проведенный анализ показал необходимость разработки интегрированного подхода, который бы объединил преимущества современных генеративных моделей с гарантиями дифференциальной приватности, что и составляет основу для дальнейшей работы по созданию практического решения.

%На основе проведенного анализа можно сделать следующие выводы о направлениях исследования:

